% s01_intro. Section file; label prefix: int. The main theorems are stated here and proved in the body.
\section{Introduction}
\label{sec:int}

\draftnote{Working title; to be settled by the authors.}

%-------------------------------------------------------------------------------
\subsection{The setting}
\label{sec:int-setting}

Let $X$ be the resolved conifold, the total space of $\cO(-1)\oplus\cO(-1)\to\bbP^1$, and let $Z=X\times\bbC^2$. The five-torus $\Tmax=(\bbC^*)^5$ acts on $Z$ with two fixed points $p_\pm$ on the zero section, with tangent weights $(\alpha_0,\alpha_1,\alpha_2,-\epsilon_1,-\epsilon_2)$ at $p_+$ and $(-\alpha_0,\alpha_0+\alpha_1,\alpha_0+\alpha_2,-\epsilon_1,-\epsilon_2)$ at $p_-$: $\alpha_0$ is the weight of the tangent line of $\bbP^1$, $\alpha_1,\alpha_2$ are the weights of the two $\cO(-1)$ summands, and $-\epsilon_1,-\epsilon_2$ those of the two trivial directions (\cref{sec:set-target}). The \emph{Calabi--Yau torus} $\TCY\subset\Tmax$ is the four-torus on which the five tangent weights sum to zero,
\[
E\coloneqq\alpha_0+\alpha_1+\alpha_2=\epsilon_1+\epsilon_2 ,
\]
and we write $\wA\coloneqq\alpha_0+\alpha_2$, $\wB\coloneqq\alpha_0+\alpha_1$. The equivariant Gromov--Witten invariants $\GW_{g,d}$ of $Z$ in genus $g$ and degree $d[\bbP^1]$ are defined, as in \cite{BS}, by virtual localisation on the proper fixed loci of the moduli spaces of stable maps to $\bbP^1$ (\cref{def:set-inv}); on $\TCY$ they are elements of the function field $\KCY=\bbQ(\alpha_0,\alpha_1,\alpha_2,\epsilon_1)$, $\epsilon_2=E-\epsilon_1$. The genus-$g$ series is
\[
\sF_g(\cQ)=\sum_{d\geq0}\GW_{g,d}\,\cQ^d ,\qquad \cQ=\re^{t},
\]
with the constant-map term $d=0$ included; the series of \cite{BS} runs over $d\geq1$, and the difference is the constant $\sF_g(0)$. The \emph{conifold point} is $\cQ=1$, that is $t=0$, where the two critical points of the superpotential governing the genus-zero theory collide with a logarithmic singularity.

The tool is the Givental--Teleman expression of the higher-genus theory as a sum over stable graphs. On $\TCY$ the point $\cQ=0$ is semisimple (because $\alpha_0\neq0$), the equivariant classes of $Z$ form a semisimple cohomological field theory of rank two, and Teleman's classification writes it as the action of an $R$-matrix on the topological field theory with vertex weights $\Delta_i^{g_v-1}$, $i\in\{1,2\}$. The two idempotents behave very differently at the conifold point: one, the \emph{singular} idempotent (colour $1$), has $\Delta_1\sim-\epsilon_1\epsilon_2E^3/t^2$, the other, the \emph{regular} idempotent (colour $2$), has $\Delta_2$ analytic and non-zero at $t=0$. We write the graph sum in a particular frame, \emph{frame~B}: the $R$-matrix is multiplied on the right by the scalar series $\hw^{-1}$, $\hw=h_{(E,-\epsilon_1,-\epsilon_2)}$, which by Mumford's formula transfers the normalised Hodge class
\[
\Hg_{\bw}=\prod_{a=1}^{3}\frac{\Lamv(\bw_a)}{\bw_a^{g_v}},\qquad \bw=(E,-\epsilon_1,-\epsilon_2),\qquad \Lamv(u)=\sum_i(-1)^i\lambda_iu^{g_v-i},
\]
onto every vertex. The resulting \emph{frame-B Hodge-vertex graph sum} $\Gsum_g$ (\cref{def:set-graphsum}) is a sum over connected stable graphs of genus $g$ without legs, over colourings of their vertices by $\{1,2\}$, and over decorations, of products of vertex weights, dilaton factors and edge factors, all built from two rows $\rr_{0,i}(z;t)$, $\rr_{1,i}(z;t)$ per idempotent. The class $\bw$ is forced: it is the unique class, up to pairs $(v,-v)$, for which the singular row takes the value $E/(E+z)$ at the conifold point (\cref{lem:set-hmult}, \cref{cor:row-F2}). The rows themselves are normalised stationary-phase expansions of the Euler-type integrand $\re^{\cW/z}\dd U/U$ at the two critical points of the superpotential $\cW$, multiplied by one constant $\hnw=h_{(\alpha_1,\wA,-\epsilon_1,-\epsilon_2)}$ fixed at $\cQ=0$.

%-------------------------------------------------------------------------------
\subsection{Statements}
\label{sec:int-statements}

The notation is that of \cref{sec:set}: the non-degeneracy condition (ND) is $\alpha_1\alpha_2\wA\wB\epsilon_1\epsilon_2E\neq0$ (\cref{def:set-ND}), a \emph{weight point with (ND)} is a point of $\mathrm{Lie}(\TCY)$ at which (ND) holds, $\alpha_0=0$ being allowed, and ``at the generic point of $\TCY$'' means ``in $\KCY$''. The frame-B rows and factors are those of \cref{def:set-rows,def:set-factors}; $\rho(t)=(\Delta_1t^2/(-\epsilon_1\epsilon_2E^3))^{-1/2}$ with $\rho(0)=1$; the \emph{leaves} are $\lf_j(\psi)=1/(E-j\psi)$, $j\geq1$, the \emph{leaf spaces} are $\Lsp_m=\mathrm{span}(\lf_1,\dots,\lf_m)$, $\qo=E-E^2\lf_1$, $\qt(\psi)=\sum_{d\geq1}\Dil_1(d)\psi^d$ is the dilaton series of the singular idempotent, and $f_{j,b}(\psi)=\coef{w^b}\Edge_{1j}(\psi,w)$ are the legs (\cref{def:set-rho}); $\tau=1-1/\cQ$, and $\delta$, $\sast$, $\Lv_2$, $\alg$ are the algebraic functions of $\tau$ of \cref{def:set-alg}, \cref{lem:set-eig} and \cref{lem:row-hess}. The first theorem is the structural statement from which everything else follows.

\begin{thm}[Theorem~A: leaf structure at generic weights]\genusstatus{AI-proposed, unchecked}
\label{thm:int-leaf}
Let $(\alpha_0,\alpha_1,\alpha_2,\epsilon_1,\epsilon_2)$ be a weight point of $\TCY$ with (ND), or indeterminates with $E=\epsilon_1+\epsilon_2$ imposed. All factors are the frame-B factors at their true idempotents. Then:
\begin{description}
\item[(F1)] $\Delta_1=-\epsilon_1\epsilon_2E^3/(t^2\rho^2)$, where $\rho=\tau/(t\alg)$ is analytic at $t=0$ with $\rho(0)=1$, so that $\Delta_1=-(\epsilon_1\epsilon_2E^3/t^2)(1-\uq t/E^2+O(t^2))$, $\uq=(E-2\alpha_1)(E-2\alpha_2)$; $\Delta_2=\epsilon_1\epsilon_2\wB\delta\sast\Lv_2$ is analytic at $t=0$ with $\Delta_2(0)=\epsilon_1\epsilon_2\alpha_1\alpha_2\wA\wB/E\neq0$ and $\Delta_2'(0)=0$.
\item[(F2)] $\rr_{0,1}(z;0)=E/(E+z)$, and $\rho\,\rr_{1,1}(z;t)=E/t+\alpha_1\alpha_2/(E+z)+O(t)$.
\item[(F3)] Every frame-B dilaton factor $\Dil_i(d)$ ($i=1,2$, $d\geq1$) and every edge factor $\Edge_{ij}(m,n)$ ($i,j\in\{1,2\}$, $m,n\geq0$) is analytic at $t=0$.
\item[(L1)] $\coef{t^0}(\rho\qt)=\qo=E-E^2\lf_1$, and $\coef{t^k}(\rho\qt)\in\Lsp_{k+1}$ for every $k\geq1$.
\item[(L2)] For $j\in\{1,2\}$, every $b\geq0$ and every $k\geq0$, $\coef{t^k}(\rho f_{j,b})\in\Lsp_{k+1}$. Moreover $\Edge_{21}(b,a)=\Edge_{12}(a,b)$, and $\coef{t^k}f_{j,b}\in\Lsp_{k+1}$ as well.
\item[(L3)] For every $k\geq0$, $\coef{t^k}(\rho^2\Edge_{11}(\psi,\psi'))\in\mathrm{span}\{\lf_i(\psi)\lf_l(\psi'):i,l\geq1,\ (i-1)+(l-1)\leq k\}$, and $\Edge_{11}(a,b)=\Edge_{11}(b,a)$. The same support bound holds for $\coef{t^k}(\rho^e\Edge_{11})$ for every integer $e$.
\end{description}
Every membership is an identity of formal power series in $\psi$ (and $\psi'$), valid at every $\psi$-order and every $t$-order; the coefficients are rational functions of the weights, regular on (ND). The bounds in (L1) and (L3) are sharp at generic weights.
\end{thm}

In words: at the conifold point the vertex weight of the singular idempotent has a double pole with an exactly known unit, every factor is regular, and the $t$-expansion of every rescaled factor at the singular idempotent is \emph{leaf-structured}: the coefficient of $t^k$ is a combination of the simple-pole functions $1/(E-j\psi)$ with $j\leq k+1$, and of their products for loops, with no constant term. \Cref{thm:int-leaf} is proved as \cref{thm:lea-main}, where the sharpness statements and a closed form of every leaf coefficient are given (\cref{thm:lea-closed}). Its proof does not use the Calabi--Yau condition: the frame-B rows depend on $(\alpha_0,\alpha_1,\alpha_2)$ only, and $\epsilon_1\epsilon_2$ enters the factors as an overall constant, so the theorem holds at every weight point of the five-torus with (ND), with $\bw=(E,-\epsilon_1,-\epsilon_2)$ and $E=\alpha_0+\alpha_1+\alpha_2$ (\cref{rmk:lea-cy}). The Calabi--Yau condition enters the main chain once, in Theorem~C below, where it makes $\bw$ a sum-zero triple whose first weight is the leaf scale $E$.

\begin{thm}[Theorem~B: the graph sum is the potential]\genusstatus{AI-proposed, unchecked}
\label{thm:int-ident}
Let $g\geq2$. At the generic point of $\TCY$, $\sF_g$ lies in $\KCY(\cQ)$, is regular at $\cQ=0$ with expansion $\sum_{d\geq0}\GW_{g,d}\cQ^d$ there, and
\[
\sF_g=\Gsum_g ,
\]
the frame-B Hodge-vertex graph sum \eqref{eq:set-Gsum} with vertex class $\Hg_{\bw}$, $\bw=(E,-\epsilon_1,-\epsilon_2)$, and the frame-B factors of \cref{def:set-rows,def:set-factors}, in which $\RB=\RZ\hw^{-1}$ and $\RZ$ has the rows $\hnw\rpt_{0,i}$, $\nw=(\alpha_1,\wA,-\epsilon_1,-\epsilon_2)$. Each coloured term $\GCont(\Gamma,\col)$ is an element of the function field $\KCY(\Ccov)$ of the double cover $\ydc^2=\Disc(\cQ)$ of $\bbP^1_\cQ$, regular over every $\cQ\notin\{0,1,\infty\}$ with $\Disc(\cQ)\neq0$.
\end{thm}

\Cref{thm:int-ident} is proved as \cref{prop:ide-main}. It is an identity in $\KCY(\cQ)$, proved at the generic point of $\TCY$ without any specialisation: the field $\KCY$ embeds in $\bbC$, Teleman's classification applies there, and the $R$-matrix descends to $\KCY\lser{\cQ}$.

\begin{thm}[Theorem~C: regularity and the polar part]\genusstatus{AI-proposed, unchecked}
\label{thm:int-reg}
Fix $g\geq2$. In frame B with $\bw=(E,-\epsilon_1,-\epsilon_2)$, at the generic point of $\TCY$ or at any weight point with (ND):
\begin{enumerate}
\item[(a)] every coloured stable graph $(\Gamma,\col)$ of genus $g$ with at least one edge is regular at $t=0$ after summing over decorations;
\item[(b)] the single-vertex graph is regular at colour $2$, and at colour $1$ its polar part is exactly
\[
\gapg\,t^{2-2g},\qquad \gapg\coloneqq(-\epsilon_1\epsilon_2E^3)^{g-1}\ZZ[g,0]=\beta_{2g-2}(\epsilon_1,\epsilon_2)=\frac{(2g-3)!\,\Pg}{\epsilon_1\epsilon_2},
\]
with no polar term of order $t^{3-2g},\dots,t^{-1}$; here $\ZZ[g,0]$ is the Hodge integral of \cref{def:set-fpot}, and
\[
\Pg=\coef{t^{2g}}\frac{\epsilon_1\epsilon_2\,t^2}{4\sinh(\epsilon_1t/2)\sinh(\epsilon_2t/2)},\qquad
\beta_n(\epsilon_1,\epsilon_2)=(n-1)!\,\coef{t^n}\frac{1}{4\sinh(\epsilon_1t/2)\sinh(\epsilon_2t/2)} ;
\]
\item[(c)] hence the expansion of $\Gsum_g$ in $\KCY\lser t$ has polar part exactly $\gapg t^{2-2g}$, and, by \cref{thm:int-ident}, so has the image of $\sF_g$ under $\cQ=\re^t$.
\end{enumerate}
\end{thm}

\Cref{thm:int-reg} is proved as \cref{cor:reg-main}. Parts (a), (b) and the first half of (c) are statements about the decoration sums as defined in \cref{def:set-graphsum}; they rest on \cref{thm:int-leaf} and on the two theorems of the companion paper \cite{LP2} cited in \cref{sec:reg-cited}, and on the string and dilaton equations (\cref{sec:reg-hodge,sec:reg-string}). The whole polar part of the graph sum sits in the single-vertex graph at the singular colour.

\begin{thm}[Theorem~D: rigidity]\genusstatus{AI-proposed, unchecked}
\label{thm:int-rig}
For $g\geq2$, in $\KCY(\cQ)$,
\[
\sF_g(\cQ)-\sF_g(0)=\frac{\Pg\,\Li_{3-2g}(\cQ)}{\epsilon_1\epsilon_2},\qquad\text{that is}\qquad \epsilon_1\epsilon_2\,\GW_{g,d}=\Pg\,d^{2g-3}\quad\text{for every }d\geq1,
\]
with $\Li_{-n}(\cQ)=\sum_{d\geq1}d^n\cQ^d$. Both sides of the second identity are polynomials in the weights on $\mathrm{Lie}(\TCY)$, so it holds at every point of $\mathrm{Lie}(\TCY)$, including $\alpha_1=\alpha_2$, $E=0$, $\epsilon_1=0$ or $\epsilon_2=0$, and $\alpha_0=0$; in particular $\epsilon_1\epsilon_2\GW_{g,d}$ depends on $(\epsilon_1,\epsilon_2)$ only, symmetrically and through $\epsilon_1^2,\epsilon_2^2$.
\end{thm}

\Cref{thm:int-rig} is proved as \cref{thm:rig-main}. In genus $0$ one has $\epsilon_1\epsilon_2\GW_{0,d}=d^{-3}$ at every weight, by the Aspinwall--Morrison formula (\cref{prop:rig-g0}), and in genus $1$ one has $\epsilon_1\epsilon_2\GW_{1,d}=\Ph{1}/d=-(\epsilon_1^2+\epsilon_2^2)/(24d)$ on $\TCY$ (\cref{lem:rig-g1}); the genus-one statement is proved by direct localisation for $d=1$, and for $d\geq2$ it uses, as an input, the formula of \cite{BS} for the resolved conifold on the subtori on which $\alpha_i=\epsilon_j$. Together:

\begin{cor}[the rigidity conjecture of \cite{BS} for the resolved conifold]\genusstatus{AI-proposed, unchecked}
\label{cor:int-conj}
For every $d\geq1$ and every $g\geq0$, $\epsilon_1\epsilon_2\,\GW_g(Z,\TCY;d[\bbP^1])=\Pg\,d^{2g-3}$, so that
\[
\sum_{g\geq0}\GW_g(Z,\TCY;d[\bbP^1])=\frac1d\Big(2\sinh\frac{d\epsilon_1}2\Big)^{-1}\Big(2\sinh\frac{d\epsilon_2}2\Big)^{-1}.
\]
Hence the rigidity conjecture of \cite{BS} holds for $X$ the resolved conifold, every non-zero class $d[\bbP^1]$, and the group $\TCY$; by contravariance it holds for every subtorus of $\TCY$ on which the invariants are defined by localisation (the normal bundle $\Nb$ moving on the fixed loci) and for every connected group whose maximal torus acts on $Z$ through such a subtorus, in particular for every Calabi--Yau subtorus of the five-torus with this property. \draftnote{BS label: conj: rigidity.}
\end{cor}

\Cref{cor:int-conj} is \cref{cor:rig-conj}, where the conjecture is stated by content. The three-torus $\Tthree=\{\alpha_0=0\}\cap\TCY$, which contains the fibrewise torus $\TF$ ($\alpha_i=\epsilon_i$) on which the zero section is fixed pointwise, is not covered by the proof of \cref{thm:int-ident} directly, because $\cQ=0$ is not semisimple there; Theorems~A and~C hold on $\Tthree$ by a rescaling of the rows to the fibrewise ones and a ring map applied to the identity of \cref{thm:int-ident} (\cref{cor:rig-3T}), and rigidity on $\Tthree$ is the restriction of \cref{thm:int-rig} by polynomiality.

%-------------------------------------------------------------------------------
\subsection{Inputs}
\label{sec:int-inputs}

The results are conditional only on the following cited theorems, each named at the place where it enters.
\begin{enumerate}
\item Stirling's series for $\log\Gamma$ with remainder \cite[5.11.1 and Section~5.11(ii)]{DLMF}, together with the Beta integral \cite[5.12.1]{DLMF} and Laplace's method, which is proved in full (\cref{lem:row-laplace}). This is the only analytic input of Theorem~A; it enters through the normalised expansion of the Beta phase (\cref{lem:row-betaphase}, \cref{prop:row-beta}).
\item Teleman's classification of semisimple cohomological field theories \cite[Theorem~(2), (1.15), Prop.~8.7(ii)]{Tel12}, with the stability range of \cite[2.17--2.19, (2.15)]{Tel12} as reported there from \cite{Loo96,MW07,Har85,Iva90}; the $R$-matrix action and its group property \cite[Defs.~2.2, 2.13, Prop.~2.14]{PPZ15}; the virtual localisation formula \cite{GP99}; the twisted-theory axioms \cite[App.~1]{CG07}; Mumford's relation and formula \cite{Mum83}, \cite[Section~1.1]{FP00}; the non-vanishing of $\lambda_g\lambda_{g-1}\lambda_{g-2}$ \cite[Theorem~4]{FP00}; and the Aspinwall--Morrison formula \cite[p.~3]{FP00}. These enter Theorem~B; the localisation formula \cite{GP99} enters also Theorem~D (\cref{lem:rig-poly}) and genus $1$ (\cref{lem:rig-g1}), Mumford's relation enters genus $1$, and the Aspinwall--Morrison formula is the genus-zero statement (\cref{prop:rig-g0}). \draftnote{Teleman is cited with the numbering of the arXiv version; the journal numbering has not been compared.}
\item The two theorems of the companion paper \cite{LP2}: the vanishing theorem for the Hodge integrals $\Xiz_h(\msM)$ with the class $\Hg_{(1,\tau,-1-\tau)}$ at leaf degree $\dg(\msM)\leq2h-3$ (\cite[Theorem~8.1, Corollaries~8.2--8.4]{LP2}, quoted as \cref{thm:reg-L}), and the value of the single vertex, $(-\epsilon_1\epsilon_2E^3)^{g-1}\ZZ[g,0]=\beta_{2g-2}(\epsilon_1,\epsilon_2)$ (\cite[Theorem~9.1]{LP2}, quoted as \cref{thm:reg-VI}). Both rest on the Mari\~no--Vafa formula \cite[(1)--(3)]{LLZ03}. They enter Theorem~C only.
\item For genus $1$ and $d\geq2$ only, the formula of \cite{BS} for the resolved conifold on the subtori on which $\alpha_i=\epsilon_j$ (\cref{lem:rig-g1}). \draftnote{BS label: prop:ResConifold, eq:resconGWref.}
\end{enumerate}
Which result uses which: Theorem~A uses (i); Theorem~C uses Theorem~A and (iii); Theorem~B uses (i), through the normalisation constant at $\cQ=0$, and (ii); Theorem~D uses Theorems~B and~C and, in the polynomiality lemma, \cite{GP99}; genus $0$ uses the Aspinwall--Morrison formula; genus $1$ uses \cite{GP99}, Mumford's relation and, for $d\geq2$, (iv). None of the arguments of this note uses topological recursion, a mirror theorem, a degeneration argument or a finite-order computation; the one input proved by degeneration is (iv), used in genus $1$ for $d\geq2$ only; the finite-order checks that accompanied the work are collected in \cref{sec:evi}, and no proof step refers to them.

%-------------------------------------------------------------------------------
\subsection{What is not claimed}
\label{sec:int-notclaimed}

\begin{enumerate}
\item Nothing is claimed off $\TCY$. Theorem~A holds at every weight point of the five-torus with (ND), but Theorem~C does not: off $\TCY$ the class $\bw$ does not sum to zero, the vanishing theorem does not apply, and the coloured graphs with an edge acquire poles at $t=0$ (\cref{sec:evi-offCY}, as evidence only).
\item The constant-map term $\sF_g(0)=\GW_{g,0}$ is not determined. It is a rational function of all five weights (\cref{rmk:rig-d0}), it depends on $\alpha$ (\cref{rmk:rig-F0}), and it is not part of the conjecture, which concerns non-zero classes. The only statement about it proved here is the flop identity $\sF_g(\infty)=\sF_g(0)$ (\cref{prop:rig-flop}).
\item The conjecture for group actions on $Z$ that are not conjugate into the five-torus is not addressed.
\item The proofs of the cited theorems of \cref{sec:int-inputs} are not reproduced. In particular the two theorems of \cite{LP2} are quoted, with the scale conversion written out (\cref{lem:reg-scale}).
\end{enumerate}

%-------------------------------------------------------------------------------
\subsection{Strategy}
\label{sec:int-strategy}

\emph{The singular row (\cref{sec:row}).} In the Euler coordinate $U=1-\tau s$, $\tau=1-1/\cQ$, the superpotential becomes the Euler phase $\phi_\tau(s)=\wA\log s+\alpha_1\log(1-s)-\alpha_2\log(1-\tau s)$, whose critical points stay separated at $t=0$ although the critical points of $\cW$ collide. The point-frame row at the singular idempotent is the normalised formal Laplace expansion of $\re^{\phi_\tau/z}(1-\tau s)^{-1}\dd s$ at the critical point $\sast$ (\cref{prop:row-wickrow}); a formal Laplace calculus (coordinate invariance, parameter derivatives with a moving expansion point, and the vanishing of total derivatives) shows that this expansion, suitably conjugated, satisfies Gauss's hypergeometric equation in $\tau$ (\cref{prop:row-gaussode}). Uniqueness of the solution gives the closed form (\cref{thm:row-closed})
\[
\rr_{0,1}(z;t)=\frac{E}{E+z}\,{}_2F_1\Big(1+\frac{\alpha_2}z,1+\frac{\wA}z;2+\frac Ez;\tau\Big)\exp\Big(-\frac{\dphi(\tau)}z\Big)\alg(\tau),
\]
in which the constant $E/(E+z)$ comes from the Laplace expansion of a Beta integral, evaluated by Stirling's series (\cref{prop:row-beta}), and $\dphi$, $\alg$ are explicit algebraic functions of $\tau$. The regularity of the $t$-coefficients in $z$ follows (\cref{cor:row-reg}, and formally again by an eikonal and transport decomposition, \cref{prop:row-transport}).

\emph{Leaves (\cref{sec:lea}).} The key formula $\rho\,\rr_{0,1}=(E/(E+z))(\tau/t)\,{}_2F_1\exp(-\dphi/z)$ (\cref{prop:lea-key}) exhibits, order by order in $t$, a rational function of $z$ with simple poles at $z=-E/j$, $j\leq k+1$, vanishing at $\infty$ for $k\geq1$: this is (L1), by partial fractions. Row~2 of the quantum differential equation is the operator form of Gauss's equation (\cref{prop:lea-row2}); self-adjointness of the quantum product gives the symplectic constant $\Kone=1$ (\cref{prop:lea-K1}) and a first-order flow for the loop $t^2\rho^2\Edge_{11}$ whose right side is a rank-one product of leaf series (\cref{thm:lea-flow}); integrating from $t=0$ gives (F3) for $\Edge_{11}$ and (L3). At the regular idempotent, Row~2 is proved by the same Laplace calculus at the second critical point (\cref{thm:lea-row2reg}), which yields the orthogonality of the two rows, the regularity of the mixed edge and (L2) (\cref{prop:lea-L2}). The value of the regular row at $t=0$ is a second application of the Beta-phase lemma (\cref{prop:lea-regrows}).

\emph{Identification (\cref{sec:ide}).} Over $\KCY$, the point $\cQ=0$ is semisimple, so Teleman's theorem gives a unique symplectic $R$-matrix; a flatness argument of the kind used for the quantum differential equation, with a stability input to kill the defect, shows that it solves the quantum differential equation (\cref{prop:ide-flat}). The rows $\hnw\rpt_{0,i}$ form a second flat unitary solution (\cref{prop:ide-rows}); two such solutions differ by constant diagonal exponentials of odd series (\cref{lem:ide-ambig}), which are fixed at $\cQ=0$ by the Hodge integrals of constant maps and the non-vanishing of $\mathrm{ch}_{2k-1}(\Hdg)$ (\cref{prop:ide-const}). The passage to frame~B is the group property of the $R$-action for the scalar $\hw$ (\cref{prop:ide-frameB}).

\emph{Regularity (\cref{sec:reg}).} At a singular vertex of genus $h\geq2$ the dilaton homogeneity of the Hodge potential (\cref{lem:reg-H0}) rescales every half-edge by one factor $\rho$ and converts $\Delta_1^{h-1}$ into $(-\epsilon_1\epsilon_2E^3)^{h-1}t^{2-2h}$ exactly, by (F1). Expanding the rescaled factors in leaves by (L1)--(L3), the coefficient of $t^{2-2h+k}$ is a combination of Hodge integrals $\Xiz_h(\msU)$ with leaf degree $\dg(\msU)\leq k$; the parts equal to $1$ are stripped by the string equation (\cref{lem:reg-string}), and the vanishing theorem kills every term with $k\leq2h-3$. A count over the vertices then shows that every coloured graph with an edge is regular, while the single vertex keeps exactly the term $(-\epsilon_1\epsilon_2E^3)^{g-1}\ZZ[g,0]t^{2-2g}$ (\cref{thm:reg-main}), whose value is the cited \cref{thm:reg-VI}.

\emph{Rigidity (\cref{sec:rig}).} By Theorem~B, $\sF_g$ is a rational function of $\cQ$ whose poles can only lie over $\cQ=1$, $\cQ=\infty$ or the two roots of the discriminant $\Disc(\cQ)$ of the critical equation (\cref{lem:rig-poles}). There is no pole on $\Disc=0$: the invariants $\epsilon_1\epsilon_2\GW_{g,d}$, $d\geq1$, are polynomials in the weights by localisation on a proper space (\cref{lem:rig-poly}), the discriminant has a root of positive $\alpha_0$-adic valuation, and an $\alpha_0$-adic evaluation of the identity $N=(1-\cQ)^n\Disc^m\sG_g$ forces $m=0$ (\cref{lem:rig-disc}). There is no pole at $\cQ=\infty$, because every factor is a power series in $1/\cQ$ there (\cref{lem:rig-inf}). At $\cQ=1$ the polar part is given by Theorem~C and equals that of $\Pg\Li_{3-2g}(\cQ)/(\epsilon_1\epsilon_2)$ (\cref{lem:rig-li}); the difference is a rational function without poles, hence a constant, which vanishes at $\cQ=0$. The excluded loci of $\mathrm{Lie}(\TCY)$ are reached by polynomial continuation.

\emph{Dependencies.} \Cref{sec:set} is definitions and genus-zero data. \Cref{sec:row} uses \cref{sec:set}; \cref{sec:lea} uses \cref{sec:set,sec:row}; \cref{sec:ide} uses \cref{sec:set} and the Laplace calculus and Beta-phase lemma of \cref{sec:row}, not Theorem~A; \cref{sec:reg} uses \cref{sec:set}, Theorem~A and the cited theorems of \cite{LP2}, and uses Theorem~B only for the second half of \cref{cor:reg-main}(c); \cref{sec:rig} uses Theorems~B and~C, the genus-zero data of \cref{sec:set}, the localisation formula \cite{GP99} (\cref{lem:rig-poly}) and, for the flop identity and the three-torus only, the Beta-phase lemma \cref{lem:row-betaphase} and \cref{cor:lea-a0}. Teleman's theorem and the other inputs of item (ii) of \cref{sec:int-inputs} enter only through \cref{sec:ide}, that is, only through the identification of the graph sum with $\sF_g$.

%-------------------------------------------------------------------------------
\subsection{Relation to the results of \texorpdfstring{\cite{BS}}{[BS]}}
\label{sec:int-BS}

The paper \cite{BS} defines the refined invariants by localisation, formulates the rigidity conjecture for a rigid class on $X\times\bbC^2$ under a Calabi--Yau group action, proves the case in which the action on $X$ is itself Calabi--Yau (there the invariants are the unrefined ones times $(\epsilon_1\epsilon_2)^{g-1}$, by Mumford's relation), and proves for the resolved conifold the closed formula of \cref{cor:int-conj} on the subtori on which $\alpha_i=\epsilon_j$, by degeneration of the base curve and the relative invariants of local curves; it also verifies the formula on $\TCY$ by computer in low degree and genus. \draftnote{BS label: defn: refined GW invariant; conj: rigidity; prop: rigidity unref limit; prop:ResConifold.} What is new here is the statement on all of $\TCY$, for every $d\geq1$ and every genus $g\geq2$, genus by genus, and the method: for $g\geq2$ no degeneration is used, and the result comes from the structure of the Hodge-vertex graph sum at the conifold point, that is, from Theorems~A, B and~C. The genus-one input of \cite{BS} at $d\geq2$ is the one place where a result of \cite{BS} is used (\cref{rmk:rig-g1input}).

%-------------------------------------------------------------------------------
\subsection{Plan of the note}
\label{sec:int-plan}

\Cref{sec:set} fixes the target, the tori, the invariants, the genus-zero data (superpotential, Frobenius algebra, idempotents, discriminant), the rows and factors of the frame-B graph sum, the leaf spaces and the Hodge-side integrals. \Cref{sec:row} computes the singular row in closed form. \Cref{sec:lea} proves Theorem~A, with the closed forms of the leaf coefficients and the case $\alpha_0=0$. \Cref{sec:ide} proves Theorem~B. \Cref{sec:reg} quotes the two theorems of \cite{LP2} and proves Theorem~C. \Cref{sec:rig} proves Theorem~D, the genus-zero and genus-one statements, \cref{cor:int-conj}, the flop identity and the three-torus corollary. \Cref{sec:evi} records the finite-order checks; nothing in the body refers to it.


\subsection*{Acknowledgements}

This work was supported by the UKRI grant UKRI2394 and by the Leverhulme Trust grant LIP-2023-007.
